\subsection{1851 Ceza Kanunnamesinin İçeriği}
\indent\indent 1851(1267) tarihli ceza kanunnamesi mukaddime, üç fasıl ve kırk üç maddeden meydana gelmektedir. Mukaddimede Tanzimat Fermanı ile can, mal, ırz ve namus emniyetini devletin garanti ettiği vurgulanmış,
fakat önceki kanunnamenin devletin garanti ettiği emniyeti yeterince sağlayamadığı belirtilmiştir. Bu sebeple, Kanunname-i Cedid bu durumları açıklığa kavuşturmak için can emniyeti, ırz ve namus emniyeti ve mal emniyeti 
olmak üzere üç ana fasıl olarak düzenlenmiştir.

Birinci fasıl, devlete karşı işlenen ve katl suçları üzerine on yedi maddeden oluşmaktadır. 1840 ceza kanunnamesindeki on madde ufak değişikliklerle(cezanın güncellenmesi, müesseselerin sorumlulukları vb.) birinci fasılda tekrar edilmiştir(Tablo~\ref{tab:1851-1840}). 
Geri kalan yedi maddenin ikisi katl suçunda taksir, teammüd ve azmettirme durumlarında verilecek cezaları düzenlenmiştir. Diğer iki madde ise maktulun varisleri tarafından katilin affı, maktulun diyar-ı aherede katli ve maktulun varisi olmaması
durumlarına açıklık getirmektedir. On beşinci madde ise kadınların katil ya da katle yardımcı olma durumlarını düzenlemektedir. Eğer katil kadın ise, erkeklere uygulanan cezanın uygulanması gerektiği belirtilmektedir. Kadın katle yardımcı olmuş ise,
ıslah oluncaya kadar hapsolunması buyrulmaktadır. Erkeler için aynı suçun cezası on dördüncü maddede düzenlenmiş ve bir seneden üç seneye kadar kürek ya da pıranga olarak belirtilmiştir.

İkinci fasıl, ırz ve namusa karşı işlenen suçlar üzerine yedi maddeden oluşmaktadır. Bu fasıl aslında 1840 Ceza Kanunnamesi'nin üçüncü faslının daha detaylandırılmış hali gibidir. 1840 Ceza Kanunnamesi'nin üçüncü faslının birinci maddesinin 
ihtiva ettiği konular ayrı maddeler halinde düzenlenmiş. Bununla birlikte sokakta sarkıntılık edenler, sarhoşlar, kumarbazlar için beşinci madde düzenlenmiş. Altıncı madde ise kız kaçırıma suçu mahalli mahkemelerce tahkik edilip, hüküm olunması
üzerinedir.

Üçüncü fasıl ise ağırlıklı olarak mal ve emlâk güvenliği üzerine işlenen suçları ve cezalarını detaylandırmıştır. Bu fasıldaki ilk on madde, 1840 Ceza Kanunnamesi'ndeki bazı cezaların güncellenmiş halidir. 1840 Ceza Kanunnamesi'nde olmayan bazı maddeler ise şöyledir:
\begin{itemize}
    \item \textbf{On Dördüncü Madde}\\
    Zahirelerini öşür olarak vermeyenlerin suçu ve cezası üzerine
    \item \textbf{On Altıncı Madde}\\
    Hasta mahbusların afları üzerine
    \item \textbf{On Yedinci Madde}\\
    Fakir suçlulardan mahsubiyetleri süresince yedirip içirecek yakınları bulunmayanların masraflarının mîri mallardan karşılanması üzerine
    \item \textbf{On Sekizinci Madde}\\
    Hizmetkarların çiftlik tohumlarını çalması üzerine
    \item \textbf{On Dokuzuncu Madde}\\
    Esnafın hileli mal satması üzerine
\end{itemize}
\begin{table}[htbp]
\centering
\begin{tabular}{ |l|l| }
\hline
\textbf{1851 Ceza Kanunnamesi} & \textbf{1840 Ceza Kanunnamesi} \\
\hline
Birinci Fasıl Birinci Madde & Birinci Fasıl Birinci Madde \\
\hline
Birinci Fasıl İkinci Madde & Birinci Fasıl Üçüncü Madde \\
\hline
Birinci Fasıl Üçüncü Madde & Birinci Fasıl Dördüncü Madde \\
\hline
Birinci Fasıl Beşinci Madde & İkinci Fasıl Birinci Madde \\
\hline
Birinci Fasıl Altıncı Madde & İkinci Fasıl İkinci Madde \\
\hline
Birinci Fasıl Yedinci Madde & İkinci Fasıl Üçüncü Madde \\
\hline
Birinci Fasıl Sekizinci Madde & İkinci Fasıl Dördüncü Madde \\
\hline
Birinci Fasıl Dokuzuncu Madde & On Birinci Fasıl Birinci Madde \\
\hline
Birinci Fasıl On Altıncı Madde & Dokuzuncu Fasıl Üçüncü Madde \\
\hline
Birinci Fasıl On Yedinci Madde & Onuncu Fasıl Birinci Madde \\
\hline
İkinci Fasıl Birinci Madde & Üçüncü Fasıl Birinci Madde \\
\hline
İkinci Fasıl İkinci Madde & Üçüncü Fasıl Birinci Madde \\
\hline
İkinci Fasıl Üçüncü Madde & Üçüncü Fasıl Birinci Madde \\
\hline
İkinci Fasıl Dördüncü Madde & Üçüncü Fasıl Birinci Madde \\
\hline
İkinci Fasıl Yedinci Madde & Üçüncü Fasıl Dördüncü Madde \\
\hline
Üçüncü Fasıl Birinci Madde & Dördüncü Fasıl Birinci Madde \\
\hline
Üçüncü Fasıl İkinci Madde & Altıncı Fasıl Birinci Madde \\
\hline
Üçüncü Fasıl Üçüncü Madde & \makecell[l]{Yedinci Fasıl Birinci Madde\\Yedinci Fasıl İkinci Madde} \\
\hline
Üçüncü Fasıl Dördüncü Madde & Yedinci Fasıl Üçüncü Madde \\
\hline
Üçüncü Fasıl Beşinci Madde & Yedinci Fasıl Dördüncü Madde \\
\hline
Üçüncü Fasıl Altıncı Madde & Sekizinci Fasıl Birinci Madde \\
\hline
Üçüncü Fasıl Yedinci Madde & Dokuzuncu Fasıl Birinci Madde \\
\hline
Üçüncü Fasıl Sekizinci Madde & \makecell[l]{Beşinci Fasıl İkinci Madde\\Beşinci Fasıl Üçüncü Madde} \\
\hline
\makecell[l]{Üçüncü Fasıl Dokuzuncu Madde\\Üçüncü Fasıl Onuncu Madde} & \makecell[l]{Beşinci Fasıl Dördüncü Madde\\Beşinci Fasıl Beşinci Madde} \\
\hline
Üçüncü Fasıl On Beşinci Madde & İkinci Fasıl Dördüncü Madde \\
\hline
\end{tabular}
\caption{1851 ve 1840 Ceza Kanunnamelerinin Benzer Maddeleri}
\label{tab:1851-1840}
\end{table}
